\section{A workhorse specification}\label{sec:app:workhorse}

This appendix fixes explicit functional forms for every primitive of Section~\ref{sec:sp:setup:primitives} and inverts the recycling margins for them. It is the specification behind the constructions of Appendix~\ref{sec:app:proofs} and the rest point of Appendix~\ref{sec:app:stat}, both of which cite the forms here rather than restating them, and it is the starting point of the quantitative implementation, \quantnote{}. Throughout, time arguments are suppressed except where trends are the point. Since $H$ and $\mu$ name the handled flow and the handling share, the CES composite below is written $Q$ and the returns exponent $\mu_F$.

\subsection{Functional forms}\label{sec:app:workhorse:forms}

\smalltitle{Preferences.}
\begin{align}
u\left(C\right) &= \frac{C^{1-\eta}}{1-\eta},\quad \eta>0,
&
v\left(P\right) &= \psi\,\frac{P^{1+\varphi}}{1+\varphi},\quad \psi>0,\ \varphi>0.
\label{eq:app:wh:preferences}
\end{align}
CRRA consumption utility with $\eta$ the elasticity of marginal utility of~\eqref{eq:sp:sum:Cgrowth}, and power disutility of pollution with $v'>0$, $v''>0$; note $v'\left(0\right)=0$, so marginal pollution damage at a clean rest point comes entirely through output.

\smalltitle{Final goods.} A CES aggregate of capital and effective materials, subject to decreasing returns and exponential pollution damage:
\begin{align}
F\left(K^Y,R,P,t\right) &= A(t)\,e^{-\kappa P}\,Q^{\mu_F}\ \text{for }R\ge\bar R,\quad 0\text{ otherwise},
&
Q &= \left[\beta\left(K^Y\right)^{1-1/\varsigma}+\left(1-\beta\right)Z^{1-1/\varsigma}\right]^{\frac{\varsigma}{\varsigma-1}},
&
Z &\equiv B(t)\,R^{e},
\label{eq:app:wh:F}
\end{align}
with $R^e=R-\bar R$, $\kappa>0$, $\beta\in\left(0,1\right)$, $\mu_F\in\left(0,1\right)$ and elasticity of substitution $\varsigma>0$ ($\varsigma=1$ is Cobb--Douglas, $Q=\left(K^Y\right)^{\beta}Z^{1-\beta}$; $\varsigma\to\infty$ is linear, $Q=\beta K^Y+\left(1-\beta\right)Z$). The returns parameter $\mu_F<1$ stands in for a fixed factor and is what makes the technology concave in $K$ alone. The substitution parameter fixes what the technology does as materials vanish, which is what the stationary benchmark of Appendix~\ref{sec:app:stat} classifies by:
\begin{center}
\begin{tabular}{lll}
\toprule
Case & Level & Margin \\
\midrule
$\varsigma\le1$ & essential: $F\left(K,0,\cdot\right)=0$ & essential: $F_R\to\infty$ \\
$1<\varsigma<\infty$ & inessential: $F^\infty\left(K,0,0\right)=A_\infty\beta^{\frac{\mu_F\varsigma}{\varsigma-1}}K^{\mu_F}$ & essential: $F_R\to\infty$ \\
$\varsigma=\infty$ & inessential: $F^\infty\left(K,0,0\right)=A_\infty\left(\beta K\right)^{\mu_F}$ & inessential: finite choke value \\
\bottomrule
\end{tabular}
\end{center}
The three rows are, respectively, the collapse case, the exhaustion-squeeze case and the choke case of Appendix~\ref{sec:app:stat}. The table is written for $\bar R=0$, and only there does it classify anything. With $\bar R>0$ the hard floor makes materials essential in level whatever $\varsigma$ --- $F=0$ below the floor by construction, and the ``level'' column above then merely describes output \emph{at} the floor --- while the margin column, read at $R\to\bar R^{+}$ rather than $R\to0$, loses its decisiveness: in all three rows the material era ends by the discrete shutdown of Section~\ref{sec:sp:opt:foc}, with material use stopping at a strictly positive flow, so an infinite marginal value at the threshold is never reached along an optimal path. The role of $\varsigma$ under $\bar R>0$ is therefore confined to output above the floor, exactly as the taxonomy of Section~\ref{sec:sp:longrun:taxonomy} says: with a floor, essentiality is not a parameter.

\smalltitle{Recycling.} An exponential approach to the yield ceiling,
\begin{align}
a\left(x\right) &= \bar a\left(1-e^{-\xi x}\right),
&
a'\left(x\right) &= \bar a\,\xi\,e^{-\xi x},
&
\alpha\left(x\right) &= \bar a\left[1-\left(1+\xi x\right)e^{-\xi x}\right],
\label{eq:app:wh:a}
\end{align}
with $\bar a\in\left(0,1\right]$ and $\xi>0$ constant --- recycling technology is held asymptotically stationary, per Assumption~\ref{ass:lr:tech}. All Section~\ref{sec:sp:setup:primitives} properties hold: $a\left(0\right)=0$, $a'>0$, $a''<0$, $a\to\bar a$, and $\alpha\in\left[0,\bar a\right)$ with $\alpha\left(0\right)=0$. The form has an exponential tail with rate $\xi$ --- the shortfall from the ceiling decays exponentially in recycling intensity, exactly: $1-a/\bar a=e^{-\xi x}$. This is not innocuous under the soft ceiling $\bar a=1$, where the tail rather than the ceiling carries the long run, and the power-tail alternative
\begin{align}
a\left(x\right) &= \bar a\,\frac{\left(\xi x\right)^{\psi}}{1+\left(\xi x\right)^{\psi}},
&
1-\frac{a}{\bar a} &= \frac{1}{1+\left(\xi x\right)^{\psi}}\sim\left(\xi x\right)^{-\psi},
\label{eq:app:wh:a-power}
\end{align}
should be carried alongside it whenever $\bar a=1$ is the case of interest, since it yields a much
slower approach to full recovery. One restriction has to travel with it. Near the origin
$a\left(x\right)\approx\bar a\left(\xi x\right)^{\psi}$, which is \emph{convex} for $\psi>1$: the
maintained assumption $a''<0$ fails there, $a'$ is non-monotone, and the marginal recycling yield
$\alpha=a-a'x$ is strictly negative for small $x$. The form therefore satisfies the
Section~\ref{sec:sp:setup:primitives} properties globally only for $\psi\le1$. For $\psi>1$ it
remains a legitimate description of the \emph{tail} --- which is all the closure criterion of
Section~\ref{sec:app:workhorse:cbgp} uses --- but two things change and must be handled explicitly:
the corner test on the recycling-capital margin is no longer decided by $a'\left(0\right)$, because
$a'$ rises before it falls, and the interior condition can have two roots, of which the relevant one
is on the decreasing branch. Under the hard ceiling the two families are interchangeable in the
region that matters and we use~\eqref{eq:app:wh:a}.

\smalltitle{Extraction and exploration.} Power costs with a linear term, so that marginal cost at zero activity is positive (a well-defined choke) and diverges as the relevant stock runs out:
\begin{align}
C^N &= c_N(t)\left(\frac{S_{\mathrm{ref}}}{S}\right)^{\mu_N}\left[\kappa_N N+\frac{N^{1+\chi_N}}{1+\chi_N}\right],
&
C^D &= c_D(t)\left(\frac{\bar X_{\max}-X_0}{\bar X_{\max}-X}\right)^{\mu_D}\left[\kappa_D D+\frac{D^{1+\chi_D}}{1+\chi_D}\right],
\label{eq:app:wh:costs}
\end{align}
with all parameters strictly positive and $S_{\mathrm{ref}}$ a normalising constant. Marginal costs are
\begin{align}
C^N_N &= c_N\left(\frac{S_{\mathrm{ref}}}{S}\right)^{\mu_N}\left(\kappa_N+N^{\chi_N}\right),
&
C^D_D &= c_D\left(\frac{\bar X_{\max}-X_0}{\bar X_{\max}-X}\right)^{\mu_D}\left(\kappa_D+D^{\chi_D}\right),
\label{eq:app:wh:mc}
\end{align}
so $C^D_D\left(0,X,t\right)\to\infty$ as $X\to\bar X_{\max}$, delivering the exhaustible-discovery property~\eqref{eq:sp:setup:exhaustible-discovery}, and $C^N_N\left(0,S,t\right)\to\infty$ as $S\to0$, which is what sustains interior extraction on squeeze paths. All other Section~\ref{sec:sp:setup:primitives} sign conditions are immediate.

\smalltitle{Waste handling and decay.}
\begin{align}
c^T\left(\varpi,t\right) &= c_T(t)\,\frac{\varpi^{1+\chi_T}}{1+\chi_T},\quad\chi_T>0,
&
\theta\left(P\right) &= \theta_{\min}+\left(\theta_0-\theta_{\min}\right)e^{-\theta_P P},
\label{eq:app:wh:waste-decay}
\end{align}
plus a collection charge $c^c(t)$ per treated tonne, so that $c^W\left(\varpi\right)=c^c\varpi+c_T\varpi^{1+\chi_T}/\left(1+\chi_T\right)$. The treatment cost satisfies $c^T\left(0\right)=c^{T\prime}\left(0\right)=0$ with $c^c+c^{T\prime}\left(1\right)=c^c+c_T$ finite, so the upper corner $\varpi=1$ is reachable, as~\eqref{eq:sp:setup:cT-finite} requires throughout. A form diverging at $\varpi=1$ would violate that assumption and, as noted in Section~\ref{sec:sp:opt:foc}, would act as a hard circularity ceiling in disguise; it is not an admissible variant here. Decay has $\theta\left(0\right)=\theta_0>\theta_{\min}>0$ and $\theta'<0$, with $\theta_0<1$ so that the per-period retention factor $1-\theta\left(P\right)$ stays positive, respecting the regeneration floor of Assumption~\ref{ass:lr:theta}; a sufficient no-tipping condition is
\begin{align}
\rho+\theta_{\min} &\ge \frac{\theta_0-\theta_{\min}}{e},
\label{eq:app:wh:notipping}
\end{align}
since $\theta+\theta'P\ge\theta_{\min}-\left(\theta_0-\theta_{\min}\right)\theta_P\max_P Pe^{-\theta_P P}$ and $\max_P Pe^{-\theta_P P}=1/\left(e\theta_P\right)$.

\smalltitle{Trends and constants.} Under Assumption~\ref{ass:lr:tech}, $A(t)=A_0e^{g_At}$ and $B(t)=B_0e^{g_Bt}$, and every remaining coefficient converges exponentially, $z(t)=z_\infty+\left(z_0-z_\infty\right)e^{-g_z t}$ for $z\in\left\{c_N,c_D,c^c,c_T\right\}$, the costs falling. In the stationary benchmark of Appendix~\ref{sec:app:stat}, $A$ and $B$ converge in the same way, rising. The accounting parameters $\omega^Y,\omega^N,\omega^D,\omega^W,\omega^C,\phi^I,d^W,\Omega^{N,S},\Omega^{D,S}$ are constants.

\subsection{The recycling margins in closed form}\label{sec:app:workhorse:margins}

Recall the material-adjusted marginal product of capital $\tilde F_K\equiv F_K\left(1-\zeta\Omega^Y\right)$ of Appendix~\ref{sec:app:proofs} and the marginal gains $\mathcal G^K$ and $\mathcal G^\varpi$ from~\eqref{eq:sp:sum:gains}. The exponential yield function makes both recycling conditions invertible.

\smalltitle{Recycling capital.} The interior condition $a'\left(x\right)\mathcal G^K=\tilde F_K$ reads $\bar a\xi e^{-\xi x}\mathcal G^K=\tilde F_K$, so
\begin{align}
x^* &= \frac{1}{\xi}\ln\frac{\bar a\xi\,\mathcal G^K}{\tilde F_K},
&
a^* &= \bar a-\frac{\tilde F_K}{\xi\,\mathcal G^K},
&
\alpha^* &= \bar a-\frac{\tilde F_K}{\xi\,\mathcal G^K}\left[1+\ln\frac{\bar a\xi\,\mathcal G^K}{\tilde F_K}\right],
\label{eq:app:wh:xstar}
\end{align}
whenever $\bar a\xi\,\mathcal G^K>\tilde F_K$, and $x^*=0$ otherwise. The middle expression is worth displaying on its own: \emph{the optimal yield is the ceiling minus the ratio of the capital cost to the (scaled) value of displaced virgin material.} Circularity approaches its technological maximum exactly as $\mathcal G^K\to\infty$, and recycling shuts down exactly when the first unit of recycling capital cannot pay for itself, $\bar a\xi\,\mathcal G^K\le\tilde F_K$.

\smalltitle{Treatment share.} The interior condition $\mathcal G^\varpi=\left[c^c+c_T\varpi^{\chi_T}\right]\left(1+\zeta\Omega^W\right)$ inverts to
\begin{align}
\varpi^* &= \min\left\{1,\;\left[\frac{\mathcal G^\varpi/\left(1+\zeta\Omega^W\right)-c^c}{c_T}\right]^{1/\chi_T}\right\}
\qquad\text{when }\mathcal G^\varpi>c^c\left(1+\zeta\Omega^W\right),
\label{eq:app:wh:varpistar}
\end{align}
and $\varpi^*=0$ otherwise: treatment shuts down when the gain cannot cover the collection charge on the first tonne.

\smalltitle{Extraction and exploration.} Writing $\Psi^{\mathrm{net}}\equiv\Psi-p^S-p^W\left(1+\Omega^{N,S}\right)$ for the value of a virgin tonne net of its rent and waste charges, the interior extraction condition inverts to
\begin{align}
N^* &= \left[\frac{\Psi^{\mathrm{net}}}{c_N\left(S_{\mathrm{ref}}/S\right)^{\mu_N}\left(1+\zeta\Omega^N\right)}-\kappa_N\right]^{1/\chi_N},
\label{eq:app:wh:Nstar}
\end{align}
when the bracket is positive and $N^*=0$ otherwise; exploration is symmetric with net value $p^S+p^X-p^W\Omega^{D,S}$ and parameters $\left(c_D,\mu_D,\kappa_D,\chi_D\right)$.

\subsection{Verification status}\label{sec:app:workhorse:status}

Checked against Section~\ref{sec:sp:setup:primitives}: all sign and boundary conditions of the forms above, including~\eqref{eq:sp:setup:exhaustible-discovery}, the yield ceiling in both cases (H) and (S), and the finite marginal treatment cost~\eqref{eq:sp:setup:cT-finite}. Derived: the closed-form margins~\eqref{eq:app:wh:xstar}--\eqref{eq:app:wh:Nstar}. The constructions that use these forms, and their status, live where they are stated: the three long-run states in Appendix~\ref{sec:app:proofs}, the dematerialized rest point in Appendix~\ref{sec:app:stat}, and the convexity conditions the forms satisfy or fail in Appendix~\ref{sec:app:suff}.

Pending. \emph{(i)}~The power-tail form~\eqref{eq:app:wh:a-power} has been carried through the closure criterion of Appendix~\ref{sec:app:proofs} but not through the closed-form margins~\eqref{eq:app:wh:xstar}, which remain written for the exponential form. \emph{(ii)}~The non-Cobb--Douglas cases: every construction in Appendix~\ref{sec:app:proofs} uses $\varsigma=1$, and while $\varsigma$ cannot change the classification when $\bar R>0$, it does change the rate algebra when $\bar R=0$.
